\chapter{Fundamentação teórica}

Este capítulo reúne os conceitos que sustentam as decisões tomadas no projeto. Inicialmente são
tratados os assistentes de voz e as tecnologias de fala; em seguida, os modelos de linguagem e o uso
de ferramentas; por fim, os fundamentos de engenharia de software que orientaram a organização do
código: arquitetura hexagonal, princípios SOLID, código limpo, métricas de complexidade, testes
automatizados e qualidade de software.

\section{Assistentes pessoais por voz}

Um assistente pessoal por voz é um sistema interativo que recebe pedidos falados, interpreta a
intenção do usuário, executa ações e devolve uma resposta, também falada. Segundo
\textcite{jurafsky2009}, sistemas de diálogo orientados a tarefas combinam tradicionalmente três
componentes: o reconhecimento automático de fala, que converte o sinal acústico em texto; a
compreensão da linguagem, que identifica a intenção e os parâmetros do pedido; e a geração e a síntese
da resposta. Entre esses componentes, um gerenciador de diálogo decide quando ouvir, quando responder
e como lidar com interrupções.

A qualidade percebida de um assistente depende fortemente da latência entre o fim da fala do usuário
e o início da resposta. \textcite{nielsen1993} associa tempos de resposta próximos de um segundo à
manutenção do fluxo de pensamento do usuário; acima de dez segundos, a atenção se perde. Em um
assistente de voz, a espera inclui a detecção do fim da fala, o reconhecimento, a geração do texto e
a síntese do áudio, o que motiva arquiteturas que processam essas etapas de forma sobreposta, em vez
de sequencial.

\section{Processamento de fala}

\subsection{Detecção de atividade de voz}

A detecção de atividade de voz (\en{Voice Activity Detection}, VAD) separa os trechos com fala dos
trechos com silêncio ou ruído. Abordagens clássicas comparam a energia do sinal com um limiar estimado a
partir do ruído de fundo, enquanto abordagens estatísticas modelam a distribuição do ruído e da fala
\cite{sohn1999}. No Robofael, a detecção é usada para decidir quando o usuário começou e terminou uma
frase. Para isso, calibra-se o limiar a partir do ruído medido na abertura do microfone e exige-se uma
pausa mínima para encerrar a captura. Ruídos curtos, como cliques, são descartados.

\subsection{Reconhecimento automático de fala}

O reconhecimento automático de fala (\en{Speech-to-Text}, STT) converte o áudio em texto. Sistemas
modernos combinam modelos acústicos e de linguagem treinados em grandes volumes de dados
\cite{jurafsky2009}. O projeto emprega dois reconhecedores com papéis complementares. O primeiro é o
Vosk, construído sobre o conjunto de ferramentas Kaldi \cite{povey2011}, que roda inteiramente no
computador e permite restringir o vocabulário a uma gramática fechada \cite{alphacephei}. O segundo é
o serviço de reconhecimento do Google, mais preciso para frases livres, usado apenas depois que o
usuário chama o assistente. A gramática restrita é especialmente útil para frases de ativação curtas,
que um reconhecedor de vocabulário aberto tende a confundir com palavras mais frequentes.

\subsection{Síntese de fala e prosódia}

A síntese de fala (\en{Text-to-Speech}, TTS) produz áudio a partir de texto. Vozes neurais atuais
permitem ajustar parâmetros prosódicos, como a frequência fundamental (o tom) e a velocidade de
elocução (o ritmo). A frequência fundamental de uma voz pode ser estimada pela autocorrelação do sinal
em janelas curtas, procurando o período que maximiza a semelhança do sinal consigo mesmo dentro da
faixa típica da voz humana, entre 60 Hz e 400 Hz \cite{jurafsky2009}. O Robofael usa essa estimativa
para comparar a voz gravada do usuário com a voz sintética e calcular o deslocamento de tom e de ritmo
que as aproxima.

\section{Modelos de linguagem de grande porte e uso de ferramentas}

Os modelos de linguagem de grande porte baseiam-se na arquitetura \en{Transformer}, cujo mecanismo de
atenção permite modelar dependências de longo alcance no texto \cite{vaswani2017}. Além de gerar texto,
esses modelos podem ser instruídos a solicitar a execução de funções externas, as chamadas ferramentas
(\en{tools}): o modelo devolve o nome da ferramenta e os argumentos em formato estruturado; a aplicação
executa a função e devolve o resultado; e o modelo continua a resposta. A API do Claude oferece esse
mecanismo com transmissão em fluxo contínuo (\en{streaming}) do texto e das entradas das ferramentas
\cite{anthropic}. Quando as entradas são transmitidas antes de estarem completas, cabe à aplicação
validá-las contra o esquema declarado antes de executar qualquer ação, porque uma entrada truncada não
pode ser tratada como válida.

Modelos menores, com alguns bilhões de parâmetros, podem ser executados no próprio computador por
servidores locais como o Ollama, que expõe uma API HTTP de conversa com suporte a ferramentas
\cite{ollama}. A execução local dispensa a conexão e não envia o texto a terceiros, mas, sem placa de
vídeo, as respostas são mais lentas, e modelos pequenos tendem a escolher ferramentas com menos
precisão. Neste trabalho, essa diferença foi tratada restringindo as ferramentas oferecidas ao modelo
local (Capítulo 6).

\section{Arquitetura de software}

\subsection{Arquitetura hexagonal}

A arquitetura hexagonal, também chamada de arquitetura de portas e adaptadores, foi proposta por
\textcite{cockburn2005} com o objetivo de isolar a lógica da aplicação dos mecanismos de entrada e
saída. A aplicação define \emph{portas}, isto é, interfaces que descrevem o que ela precisa do mundo
externo ou o que oferece a ele, e cada tecnologia concreta é conectada por meio de um
\emph{adaptador} que implementa a porta. Com isso, banco de dados, interface gráfica, serviços de
voz e modelos de linguagem tornam-se substituíveis, e a lógica central pode ser testada com dublês
simples. \textcite{martin2019} generaliza a ideia na chamada regra de dependência: as dependências de
código-fonte devem apontar sempre para dentro, das camadas de detalhe para as camadas de política.

\subsection{Domínio, aplicação e raiz de composição}

\textcite{evans2004} propõe que as regras de negócio sejam expressas em um modelo de domínio
independente de infraestrutura, com entidades, objetos de valor e serviços de domínio. Os casos de uso
que coordenam esse modelo formam a camada de aplicação. A ligação entre abstrações e implementações
concretas é feita por injeção de dependências \cite{fowler2004}, idealmente concentrada em um único
ponto do programa, a raiz de composição, que monta o grafo de objetos na inicialização.

\subsection{Padrões de projeto empregados}

Diversos padrões catalogados por \textcite{gamma1994} aparecem na solução: o \en{Observer}, na forma
de um barramento de eventos que desacopla produtores e consumidores; o \en{State}, em uma máquina de
estados explícita que controla o ciclo de conversa; o \en{Strategy} e a \en{Chain of Responsibility},
nas cadeias de interpretadores de intenções e de períodos; e o \en{Adapter}, nos componentes que
encapsulam bibliotecas externas.

\subsection{Funções de aptidão arquitetural}

\textcite{ford2017} propõem tratar características arquiteturais desejadas como funções de aptidão
(\en{fitness functions}): verificações objetivas e, preferencialmente, automatizadas, executadas
continuamente para indicar se a arquitetura ainda atende ao que se espera dela. Um exemplo é um teste
que inspeciona as dependências entre módulos e falha quando uma camada interna passa a depender de uma
externa. Essa ideia fundamenta os testes de arquitetura descritos no Capítulo 7.

\section{Princípios SOLID}

O acrônimo SOLID reúne cinco princípios de projeto orientado a objetos sistematizados por
\textcite{martin2003}:

\begin{enumerate}[label=\alph*)]
  \item \textbf{princípio da responsabilidade única} (\en{Single Responsibility Principle}, SRP): um
        módulo deve ter um, e apenas um, motivo para mudar; em outras palavras, deve responder a um único
        ator ou preocupação \cite{martin2019};
  \item \textbf{princípio aberto/fechado} (\en{Open-Closed Principle}, OCP): entidades de software
        devem estar abertas para extensão e fechadas para modificação, formulação atribuída a
        \textcite{meyer1988}; acrescentar um comportamento não deve exigir a alteração de código que já
        funciona;
  \item \textbf{princípio da substituição de Liskov} (\en{Liskov Substitution Principle}, LSP):
        objetos de um subtipo devem poder substituir objetos do tipo base sem alterar as propriedades
        desejáveis do programa \cite{liskov1994};
  \item \textbf{princípio da segregação de interfaces} (\en{Interface Segregation Principle}, ISP):
        clientes não devem ser forçados a depender de métodos que não usam, o que favorece interfaces
        pequenas e específicas;
  \item \textbf{princípio da inversão de dependência} (\en{Dependency Inversion Principle}, DIP):
        módulos de alto nível não devem depender de módulos de baixo nível; ambos devem depender de
        abstrações.
\end{enumerate}

Na linguagem Python, as abstrações exigidas pelo DIP e pelo ISP podem ser expressas por protocolos
(\cod{typing.Protocol}), que definem contratos verificados estruturalmente pelo verificador de tipos,
sem exigir herança explícita \cite{pep544}. Essa característica é explorada em todo o projeto.

\section{Código limpo e refatoração}

\textcite{martin2009} caracteriza o código limpo como aquele que é fácil de ler e de modificar por
outras pessoas: nomes que revelam a intenção, funções pequenas que fazem uma única coisa, poucos
argumentos, ausência de duplicação e de efeitos colaterais ocultos, e tratamento de erros que não
obscurece a lógica. O princípio da não repetição (\en{Don't Repeat Yourself}) de
\textcite{hunt1999} complementa essas práticas. A refatoração, segundo \textcite{fowler2018}, é a
alteração da estrutura interna do software sem mudar seu comportamento observável, feita em pequenos
passos e protegida por testes automatizados. Entre as técnicas aplicadas neste trabalho estão a
extração de classe, a extração de função, a substituição de condicionais por polimorfismo ou por
tabelas de despacho e a introdução de objetos-parâmetro.

\section{Métricas de complexidade}

A complexidade ciclomática, proposta por \textcite{mccabe1976}, mede o número de caminhos linearmente
independentes no grafo de controle de uma função. Ela corresponde, de forma aproximada, ao número de
decisões mais um. Valores altos indicam funções difíceis de entender e de testar por completo, e o
limite de 10 sugerido pelo próprio autor é amplamente adotado como teto prático. Neste trabalho, a
métrica foi usada tanto para localizar os pontos que precisavam de refatoração quanto como regra
permanente do analisador estático.

\section{Testes automatizados}

O desenvolvimento orientado a testes, descrito por \textcite{beck2003}, popularizou a prática de
escrever testes automatizados pequenos e rápidos que documentam o comportamento esperado.
\textcite{cohn2009} propõe a pirâmide de testes: muitos testes unitários na base, uma quantidade menor
de testes de integração e poucos testes de ponta a ponta, que são mais lentos e frágeis.
\textcite{meszaros2007} cataloga os dublês de teste, como \en{stubs}, \en{fakes} e \en{spies}, que
substituem colaboradores reais para isolar o comportamento sob teste. A combinação de arquitetura
hexagonal com dublês permite testar a lógica de conversa sem microfone, sem alto-falante e sem rede.

\section{Qualidade de software}

A norma ISO/IEC 25010 define um modelo de qualidade de produto com características como adequação
funcional, eficiência de desempenho, compatibilidade, usabilidade, confiabilidade, segurança,
manutenibilidade e portabilidade \cite{iso25010}. Esse modelo foi usado como referência para organizar
os requisitos não funcionais do Capítulo 4. Na engenharia de requisitos, \textcite{sommerville2018} e
\textcite{pressman2016} distinguem os requisitos funcionais, que descrevem o que o sistema deve fazer,
dos requisitos não funcionais, que restringem como ele o faz, distinção também adotada neste
trabalho.

\section{Integridade de dados e históricos auditáveis}

Um histórico é auditável quando permite reconstruir o que aconteceu e detectar alterações indevidas.
Duas técnicas complementares sustentam essa propriedade no projeto. A primeira é o armazenamento
somente de acréscimo: os registros nunca são alterados nem removidos, e restrições do próprio banco de
dados recusam essas operações. A segunda é o encadeamento de resumos criptográficos: cada registro
guarda o resumo SHA-256 \cite{nist2015} do anterior, de modo que qualquer modificação, remoção ou
inserção indevida quebra a cadeia a partir do ponto alterado. O banco de dados SQLite oferece
transações atômicas e gatilhos (\en{triggers}) adequados para implementar essas garantias em um único
arquivo local \cite{sqlite}.
